% ============================================================
%  Глава 1. Экзотические деривативы и их применение на практике
% ============================================================
\chapter{Экзотические деривативы и их применение на практике}
\label{ch:exotic}

\section{Обзор экзотических деривативов}
\label{sec:exotic-overview}

Под экзотическим деривативом понимают производный финансовый инструмент,
выплата по которому отличается от выплаты стандартного европейского или
американского опциона. Чаще всего отличие выражается в зависимости от
\emph{траектории} цены базового актива на отрезке жизни контракта,
от \emph{нескольких} базовых активов одновременно, либо в нестандартной
структуре выплат (барьеры, ловушки, цифровая компонента и пр.). Ниже мы
кратко рассматриваем четыре класса инструментов, относящихся к этой
работе.

\subsection{Азиатские опционы}
\label{sec:asian}

Азиатский опцион~--- опцион, выплата по которому зависит не от значения
базового актива в момент исполнения, а от \emph{среднего} значения цены
на некотором отрезке времени. Различают арифметические и геометрические
средние, фиксированный и плавающий страйк. В арифметическом
случае выплата по call-опциону имеет вид
\begin{equation}
\label{eq:asian-payoff}
  \text{Payoff}^{\text{Asian}}_{\text{call}}
  \;=\; \left( \frac{1}{T}\int_0^T S_t\,\dd t \,-\, K \right)^{+}.
\end{equation}

Усреднение существенно снижает чувствительность к манипуляциям и
выбросам, что делает азиатские опционы естественным инструментом для
рынков с дискретной/нерегулярной ликвидностью.

\subsection{Барьерные опционы}
\label{sec:barrier}

Барьерные опционы~--- инструменты, выплата по которым появляется
(knock-in) или, наоборот, аннулируется (knock-out) при пересечении
ценой базового актива заданного уровня (барьера) хотя бы один раз на
интервале $[0,T]$. Барьеры могут быть верхними (up) или нижними (down),
дискретными или непрерывными. Экономически барьерные опционы дешевле
обычных за счёт ограничения области выплат и используются как
инструменты <<точечного>> хеджирования сценариев.

\subsection{Радужные (rainbow) опционы}
\label{sec:rainbow}

Радужные опционы~--- класс мультиактивных опционов, выплата по которым
зависит от поведения нескольких базовых активов $S^{(1)},\dots,S^{(n)}$.
Наиболее распространены варианты:

\begin{itemize}
  \item \textbf{Best-of}: выплата вида
        $\left(\max\bigl(S^{(1)}_T,\dots,S^{(n)}_T\bigr) - K\right)^{+}$.
  \item \textbf{Worst-of}: выплата вида
        $\left(\min\bigl(S^{(1)}_T,\dots,S^{(n)}_T\bigr) - K\right)^{+}$.
  \item \textbf{Spread}: выплата вида
        $\left(S^{(i)}_T - S^{(j)}_T - K\right)^{+}$.
  \item \textbf{Basket}: выплата по корзине $\sum_i w_i S^{(i)}_T$.
\end{itemize}

Между этими видами существуют простые алгебраические соотношения,
позволяющие выражать одни через другие. Например, для двух активов
выполнено
\begin{align}
  \max(S^{(1)},S^{(2)}) &= S^{(1)} + \bigl(S^{(2)}-S^{(1)}\bigr)^{+},\\
  \min(S^{(1)},S^{(2)}) &= S^{(1)} - \bigl(S^{(1)}-S^{(2)}\bigr)^{+},\\
  \max(S^{(1)},S^{(2)}) + \min(S^{(1)},S^{(2)}) &= S^{(1)}+S^{(2)}.
\end{align}
Эти соотношения мы используем далее в~\cref{sec:rainbow-decomp} при
декомпозиции цены rainbow-опциона.

\subsection{Корреляционные свопы}
\label{sec:corr-swap-intro}

Корреляционный своп~--- внебиржевой инструмент, выплата по которому
определяется разностью между реализованной и фиксированной корреляцией
заданной корзины активов. Подробное обсуждение приводится в
\cref{ch:corr}; здесь мы фиксируем лишь общую формулу выплаты
\begin{equation}
  \text{Payoff}^{\text{CorrSwap}}
  \;=\; N\bigl(\rho_{\text{real}} - K_{\text{fixed}}\bigr),
\end{equation}
где $N$~--- номинал, $K_{\text{fixed}}$~--- фиксированный страйк,
а $\rho_{\text{real}}$~--- реализованная средняя корреляция.

\section{Применение экзотических деривативов на практике}
\label{sec:exotic-practice}

Азиатские опционы широко используются для хеджирования экспозиции к
средним ценам сырьевых активов, прежде всего нефти. Соответствующая
литература описывает применение арифметических азиатских call-опционов
для защиты от роста средней цены закупаемого сырья на горизонте квартала
или года~\cite{kemna1990, leoni2014commodity}. \todoblock{Добавить
дополнительные ссылки по азиатским опционам на рынках электроэнергии и
зерна; уточнить статью по нефтехеджу.}

Барьерные опционы применяются для <<точечного>> хеджирования резких
движений в условиях кризиса; rainbow-опционы~--- в структурном
финансировании и индексных продуктах; корреляционные свопы~--- для
торговли структурой корреляций и хеджирования портфелей от системных
скачков в коррелированности активов (см. \cref{sec:corr-hedge-bubble}).

\section{Методы оценки и их сложности}
\label{sec:exotic-methods}

Базовая идея, проходящая через всю работу, состоит в следующем:
\emph{цена дериватива равна стоимости его хеджа}~\cite{shreve2004stoch}.
В модели Блэка--Шоулза для ванильных опционов это утверждение приобретает
вид замкнутой формулы, выводящейся через построение реплицирующего
портфеля. Для ряда экзотических инструментов известны замкнутые или
полу-замкнутые формулы, но в ряде важных случаев они отсутствуют.

\begin{itemize}
  \item \textbf{Ванильные опционы (call/put).} Существует замкнутая
        формула Блэка--Шоулза; цена~--- результат построения
        реплицирующего портфеля с непрерывной ребалансировкой.
  \item \textbf{Корреляционные свопы.} В целом ряде разумных моделей
        динамики корреляции существует замкнутая формула справедливого
        страйка, см.~\cref{sec:corr-models}.
  \item \textbf{Геометрические азиатские опционы.} Существует замкнутая
        формула; для арифметических~--- лишь приближения и Монте-Карло.
  \item \textbf{Rainbow-опционы.} В общем случае замкнутых формул нет;
        имеются полу-аналитические выражения для двух активов
        (\cite{stulz1982options, johnson1987options}) и приближения
        для корзин (Vorst, Levy и др.).
  \item \textbf{Барьерные опционы.} Замкнутая формула существует для
        непрерывного мониторинга в модели Блэка--Шоулза; для
        дискретного мониторинга и стохастической волатильности
        используются Монте-Карло и численные методы.
\end{itemize}

Отсутствие удобной замкнутой формулы для rainbow-опционов на корзинах
размерности $n\ge 2$ и для арифметических азиатских опционов делает
естественным использование (i)~численных и (ii)~машинно-обучаемых
методов оценки, что и определяет дальнейшую структуру работы.

\todoblock{При необходимости расширить обзор --- добавить ссылки на
работы по \emph{control variates} для Монте-Карло, на работу
\cite{glasserman2004monte} и на современные обзоры (последние 5 лет).}
